\paragraph{Nonuniformity is unavoidable if the conductor is left unrestricted.}
The counting part of A49 can always be arranged after allowing
arbitrary odd conductors. Precisely, fix a smooth projective
geometrically connected genus-two curve $C/k$ over a finite field
$k=\mathbf F_q$, a point $P\in C(k)$ and a finite set $S$ of primes containing
$2$, the characteristic $p$, and every prime dividing the order
of $\kappa_P$ in $J_C(k)$.
There are arbitrarily large extensions $k_j/k$ and quotients
\[
 H_j=J_C(k_j)_{S^c}
\]
which kill $\kappa_P$ and for which the augmentation element
$F_{H_j}$ of (OQ.3) is nonuniform.
Consequently there is a constant connected $P$-split cyclic cover
over $k_j$, of odd conductor supported outside $S$, with at least
one faithful character $L$-value nonzero modulo $p$.
In fact its entire signed $p$-orbit is nonzero. No assertion that
all faithful characters are nonzero is made.

\begin{proof}
Use the explicit integer $M$ of (PE.3) in A42, so that
$J_C(\mathbf F_{q^j})_S=J_C(k)_S$ whenever $\gcd(j,M)=1$.
Put $B=\#J_C(k)_S$ and $Q=q^j$ on this infinite set of degrees;
the progression $j\equiv1\pmod M$ already supplies an infinite set.
The projection to $H_j$ has kernel of the fixed size $B$ and
kills $\kappa_P$. The Weil bounds \cite[Theorem~19.1]{MilneAV} give
\[
 |H_j|\ge\frac{(\sqrt Q-1)^4}{B},\qquad
 Q+1-4\sqrt Q\le N_j:=\#C(\mathbf F_Q)
                         \le Q+1+4\sqrt Q.
\]
For all sufficiently large admissible $j$, these imply
$|H_j|>N_j$ and $N_j>2(B-1)+1$.

For every nonzero Jacobian class there are at most two ordered
representations as a difference of two distinct geometric points
of a genus-two curve; this is the difference bound used in recovered
successor A.10 and A.35 \cite[A.10, A.35]{ASrun}.
Hence at most $2(B-1)$ rational points
share an $H_j$-coset with another rational point: the total number
of ordered distinct pairs in the same coset is at most that bound.
There are therefore at least two singleton cosets. At most one
is zero, so a nonzero singleton coset has coefficient one in
$F_{H_j}$. Since $|H_j|>N_j$ and $P$ occupies zero, there is
also an empty nonzero coset, whose coefficient is zero.
Thus $F_{H_j}$ is not a constant multiple of the sum of all
group elements, in every characteristic and with no parity
condition on $N_j$.

Fourier inversion now gives a nontrivial character with nonzero
value. Its cyclic image has order supported outside $S$.
The original-field Lang construction and character formula of
A49 identify this value with the corresponding faithful
$L_P(\chi,1)$ modulo a fixed prime above $p$.
Since the quotient kills $\kappa_P$, $p$-power and inversion
preserve nonvanishing. This proves precisely the signed-orbit
assertion. Iterating after adjoining the primes in the selected
conductor to $S$ also gives pairwise coprime conductors, without
making those conductors prime or signed-primitive.
\end{proof}

This combines the already checked A42 primary-stabilization
argument with the recovered difference bound and Fourier identity;
it is not a new claim about Lang construction or torsion density.
It answers the nonuniformity-only supply question positively.
It does not give the unit-group hypothesis of A49: an exponent
or selected conductor can have several signed $p$-orbits, and
the other orbits can still contain vanishing factors. The exact
remaining step is to control an entire faithful conductor,
or to bound the rational primitive exponent as in A47.
The statement therefore does not close the existing universal
geometric-witness row or automatically count as main-gate movement.
